% Laget automatisk av elevark/lag_elevark.py fra sim/dobbeltspalte.html. Endringer her blir overskrevet.
\documentclass[11pt]{article}
\usepackage{elevark}
\begin{document}
\elevarkhode{Fysikk 2 · Kvantefysikk}{Dobbeltspalte med én partikkel om gangen}{https://fysikklab.no/sim/dobbeltspalte.html}{Simuleringen}

Send partikler én og én mot en vegg med to smale spalter. Hver partikkel treffer skjermen i ett bestemt punkt, som en partikkel. Men etter tusenvis av treff dukker det opp et stripemønster, som fra bølger. Hva skjer med klinkekuler, og hva skjer hvis vi sjekker hvilken spalte hver partikkel går gjennom?

\section*{Oppgaver}

\begin{oppgaveliste}
\oppgave \fase{gjett}{Gjett}Vi sender klinkekuler gjennom to spalter. Hvordan ser mønsteret på skjermen ut? Hva tror du skjer med fotoner som sendes én og én?
\svar
\oppgave \fase{test}{Test}Velg klinkekuler og trykk Start. Bytt så til fotoner med 10 partikler per sekund. Kan du se et mønster etter 50 treff? Etter 2000?
\svar
\oppgave \fase{forklar}{Forklar}Hvert foton treffer ett punkt. Hvordan kan et enkelt foton «vite» om stripene? Hva er det som har bølgeegenskaper?
\svar
\oppgave \fase{test}{Regn}Regn ut avstanden mellom stripene for lys med $λ = 550\,\mathrm{nm}$ og $d = 0{,}25\,\mathrm{mm}$. Sjekk med simuleringen. Hva skjer når du øker $d$ eller $λ$?
\svar
\oppgave \fase{test}{Regn}Velg elektroner med akselerasjonsspenning 100~V. Regn ut de Broglie-bølgelengden. Hva skjer med stripene når spenningen økes til 400~V?
\svar
\oppgave \fase{test}{Undersøk}Velg «To spalter, mål hvilken spalte». Hva skjer med mønsteret? Sammenlign med én spalte.
\svar
\oppgave \fase{utfordring}{Utfordring}Hvorfor ser vi aldri interferens med klinkekuler? Regn ut de Broglie-bølgelengden til en klinkekule på 5~g som triller med 1~m/s.
\svar
\end{oppgaveliste}

\end{document}
